% =====================================================================================
%  Expanded body sections. Included by paper.tex.
%  Every numeric claim traces to an artifact under ../code/. Nothing is typed from memory.
% =====================================================================================

% ---------------------------------------------------------------------- INTRODUCTION
\section{Introduction}

An incident arrives at a service desk. Someone must decide, within seconds, which team
should handle it. That decision is expensive to reverse: in the log studied here, incidents
reassigned at least once take \num{8.96} times longer to resolve at the median, \num{14.39}
hours against \num{1.61}, and \SI{41.1}{\percent} of incidents are reassigned at least once.
Getting the first assignment right is therefore worth roughly a ninefold reduction in median
resolution time, and the existing process fails to do so on two cases in five.

Predictive process monitoring offers to inform that decision and reports encouraging numbers
for doing so \citep{teinemaa2019outcome,ceravolo2024predictive}. The difficulty is
structural. A model trained on completed cases has access to the whole case. Attributes such
as the number of reassignments, the count of activity events, or the number of distinct
groups that touched an incident are all recorded in the log and none is knowable at the
moment of assignment. Several are near-deterministic proxies for handle time, which is
exactly the quantity the service level label thresholds. A model built from them will score
well and cannot be deployed, because the information it depends on does not exist when the
decision it informs is taken.

\citet{kaufman2012leakage} formalised this as a violation of legitimacy: a feature is
admissible only if it was observable before the target was determined.
\citet{kapoor2023leakage} document the consequences across 17 scientific fields and 294
papers. The pattern recurs wherever operational logs are modelled, in connectome-based
neuroimaging \citep{rosenblatt2024leakage}, in clinical prediction
\citep{davis2024framework,ramadan2025diagnostic}, in software build prediction
\citep{mishra2026taxonomy}, in longitudinal education modelling
\citep{fauszt2026timeconsistent}, and in financial text \citep{xue2026temporal}.

Process mining has begun to address it. \citet{weytjens2022creating} show that constructing
benchmark splits without accounting for case overlap leaks information across the boundary.
\citet{abb2024discussion} measure example leakage above \SI{80}{\percent} across standard
BPIC logs and find that a naive baseline matches a state-of-the-art deep model once that
leakage is present, a result \citet{weytjens2026david} extend to transformer and
large-language-model baselines.

What is missing is a measurement, on one log under one protocol, of both what the discipline
costs and what it buys, reported together. A cost figure alone invites the conclusion that
admissibility is not worth paying for. A benefit figure alone is unfalsifiable. This paper
reports both.

\subsection{Research questions}

\begin{description}
\item[RQ1] How much discrimination is lost when every attribute must be observable at the
moment of assignment and every group-level encoding must be estimated from past cases only?
\item[RQ2] Under that constraint, does adding attributes beyond the small sets used in prior
work still improve discrimination, and by how much?
\item[RQ3] Does choosing an admissible modelling target, rather than the convenient one,
cost accuracy or gain it?
\item[RQ4] Does a priority-aware adaptive multi-objective assignment policy improve on
standard NSGA-II under the same admissibility constraint?
\end{description}

\subsection{Contributions}

\begin{enumerate}
\item A decision-time admissibility protocol stated as two rules and implemented as
executable constraints rather than as documentation, with its cost measured at \num{0.0792}
AUC on BPI Challenge 2014, which is \SI{25.0}{\percent} of all above-chance signal (RQ1).

\item An admissible feature ladder rising from \num{0.5100} at three attributes to
\num{0.8270} at seventeen, a gain of \num{0.2560} AUC obtained without any inadmissible
attribute, reproduced across two split ratios to within \num{0.0053} (RQ2).

\item Two decision-aware scoring formulations. An assignment score aggregating seven criteria
by a weighted power mean whose exponent parameterises compensation between criteria, and a
service level risk score obtained by modelling resolution as a discrete-time hazard and
deriving breach as a survival ratio. The second reaches \num{0.8373} AUC at intake against
\num{0.7399} for a direct classifier, while being structurally unable to observe its own
label (RQ3).

\item A multi-objective assignment study in which the proposed operator set improves on its own
base algorithm on four of five metrics but is outperformed by an off-the-shelf many-objective
control, with an ablation attributing essentially the whole gain to one of its three components.
Two reproducibility defects were found in our own harness and corrected before any number in
this paper was produced. The uncorrected harness reversed this conclusion (RQ4).
\end{enumerate}

% ---------------------------------------------------------------------- RELATED WORK
\section{Related work}

\subsection{Leakage, admissibility and target encoding}

\citet{kaufman2012leakage} give the canonical formulation and the learn-predict separation
principle. \citet{kapoor2023leakage} provide an eight-type taxonomy and connect leakage to
reproducibility failure at scale. Domain treatments converge on the same shape: an
accuracy collapse once only pre-decision features are admitted, with
\citet{mishra2026taxonomy} reporting a \num{15.07} point drop on one platform after removing
leaky features and arguing that prior studies reporting 95 to 99 per cent accuracy used
contaminated attributes.

Target encoding deserves separate attention because it is the second of our two rules.
\citet{pargent2022regularized} show that regularised target encoding outperforms traditional
categorical handling, and the regularisation they describe is exactly what makes the encoding
safe only when it is refitted inside each resampling fold. An encoding fitted once over the
full dataset before splitting contaminates every subsequent evaluation and includes each
row's own label in its own predictor.

Within process mining, \citet{weytjens2022creating} address leakage in benchmark
construction and show that a plain temporal split is insufficient because training cases
overlap the test period. \citet{abb2024discussion} quantify example leakage on the standard
BPIC logs. Our protocol extends this from the split axis to the feature axis and the time
axis at once, and measures the cost rather than only prescribing the constraint.

\subsection{Predictive and prescriptive process monitoring}

\citet{teinemaa2019outcome} provide the standard outcome-prediction benchmark across 24
real-world datasets, complemented by remaining-time benchmarks \citep{verenich2019survey} and
deep-learning comparisons \citep{ramamaneiro2023learning,kratsch2021machine}. We note
explicitly that \textbf{BPI Challenge 2014 is not in the Teinemaa dataset list}: that
benchmark states it discarded the logs from 2013 and 2014 because its selection criterion
required both case and event attributes. Its reported ranges are therefore used here as
context and never as a head-to-head comparison. We make no stronger claim than that. Absence
from one benchmark is not evidence that nobody has published on this log, and we searched for
a same-task result to compare against without finding one. \citet{fertig2026revisiting} revisit the field under foundation models and find
that tabular gradient boosting remains competitive, which is why this study uses it.

Prescriptive process monitoring asks what to do rather than what will happen
\citep{kubrak2022prescriptive}. Recent work addresses resource constraints explicitly
\citep{shoush2025prescriptive} and the organisational question of whether practitioners can
act on the output \citep{kubrak2026supporting}. The framing matters here because an
admissibility constraint is only meaningful relative to a decision, and a decision is only
meaningful if someone can take it at the time the model fires.

\subsection{IT service management and this log}

BPI Challenge 2014 \citep{dongen2014challenge} has been analysed for workload impact
\citep{dees2014predictive} and with fuzzy and heuristic mining
\citep{andreswari2023analysis}, and the BPI challenge series has been surveyed
\citep{lopes2019survey}. Adjacent operational problems include incident-management
prediction \citep{sarnovsk2018predictive}, ticket reassignment prediction
\citep{schad2022predicting}, assignment-group classification \citep{gassel2024evaluating},
service level breach prediction \citep{gouryraj2021service,swain2023factors}, and escalation
with large language models \citep{liu2025tickit}. Ticket assignment has also been posed as a
multi-objective problem \citep{arain2026novel}, which is the closest prior framing to our
Section on assignment.

\subsection{Multi-objective assignment and its known limits}

NSGA-II \citep{deb2002fast} remains the reference algorithm, with NSGA-III
\citep{debjain2014evolutionary} introduced for many-objective problems where dominance-based
selection degrades. That degradation is well characterised and is still an active theoretical
subject \citep{zheng2024runtime,doerr2026difficulties,opris2025first,alghouass2025proven},
and empirical comparisons on applied problems report mixed outcomes
\citep{santos2024performance,wolschick2024evaluating,ishibuchi2025search,wang2025nsgaii,li2025learning}.
Our five objectives place this study in exactly that regime, so a result in which NSGA-III
outperforms an NSGA-II variant is consistent with theory rather than surprising.

\citet{wang2025rethinking} argue that a substantial share of claimed novelty in metaheuristic
research does not survive controlled comparison. That argument is directly relevant to the
result we report in Section~\ref{sec:assignment-results}, and it is one reason we report the
ablation rather than only the aggregate.

\subsection{Positioning against the closest prior work}
\label{sec:positioning}

Table~\ref{tab:related} places this study against the work nearest to it on each of its three
axes: leakage measurement, service level prediction on operational logs, and multi-objective
assignment.

A word on what the table deliberately does not contain. It carries a reported-performance column
only where a figure is comparable to ours, and marks the rest as not comparable rather than
printing a number that would invite a false ranking. Accuracy figures across different logs,
different labels and different decision points are not on a common scale, and assembling them
into a league table is the specific practice that produced the leakage literature this paper
builds on. \textbf{BPI Challenge 2014 is not in the \citet{teinemaa2019outcome} dataset list},
which excludes the 2013 and 2014 logs by its own stated selection criterion, so the ranges
reported there are context and not a head-to-head comparison. That is a statement about one
benchmark, not about the literature: \citet{dees2014predictive} do publish on this log, on a
different task and without a decision-time constraint.

\begin{table}[htbp]
\centering
\caption{This study against the closest prior work. A dash in the performance column means the
figure exists in the cited work but is not on a scale comparable with ours, because the task,
the log or the decision point differs.}
\label{tab:related}
\small
%% GENERATED by code/figures/make_tables.py. Do not edit by hand.
%% The caption, label and float live in the manuscript; only the numbers are here.
\begin{tabular}{L{3.15cm} L{3.0cm} L{2.35cm} L{2.1cm} L{3.35cm}}
\toprule
Study & Approach & Data & Reported & Limitation for this question \\
\midrule
\citet{kaufman2012leakage} & Formalises leakage and the learn-predict separation & Several & --- & States the condition; does not measure what enforcing it costs \\
\citet{kapoor2023leakage} & Eight-type leakage taxonomy across 17 fields & 294 papers & --- & Survey of consequences, not a measurement on one log \\
\citet{abb2024discussion} & Measures example leakage in process mining & Standard BPIC logs & $>\SI{80}{\percent}$ example leakage & Measures the split axis; does not price the feature axis \\
\citet{mishra2026taxonomy} & Removes leaky features from build prediction & CI build logs & $-15.07$ accuracy points & Different domain; no benefit side reported \\
\citet{weytjens2022creating} & Leak-free benchmark construction & Process logs & --- & Prescribes the constraint rather than costing it \\
\citet{teinemaa2019outcome} & Outcome-prediction benchmark & 24 event logs & --- & BPI 2014 not among the datasets \\
\citet{dees2014predictive} & Workload impact analysis on this log & BPI 2014 & --- & Descriptive; no decision-time constraint \\
\citet{gouryraj2021service} & Service level breach prediction & ITSM tickets & --- & No admissibility constraint stated \\
\citet{arain2026novel} & Multi-objective ticket assignment & ITSM & --- & No leakage control on the objective models \\
\citet{debjain2014evolutionary} & NSGA-III for many-objective problems & Benchmarks & --- & Not applied to assignment under admissibility \\
\midrule
\textbf{This work, unconstrained} & Reproduction of the default practice & BPI 2014, \num{46605} & 0.9050 AUC & Reported only to size the correction \\
\textbf{This work, admissible} & Rules 1 and 2 enforced & BPI 2014, \num{46605} & \textbf{0.8258} AUC & The deployable figure \\
\textbf{This work, chronological} & Admissible, evaluated forward in time & BPI 2014, \num{46605} & 0.8074 AUC & The figure under drift \\
\bottomrule
\end{tabular}


\end{table}

What is missing from every row above is a measurement, on one log under one protocol, of both
what the discipline costs and what it buys, reported together. That is the gap this paper fills.

\subsection{Aggregation, weighting and calibration}

The weighted power mean is a compensative connective in the sense of
\citet{dyckhoff1984generalized}, generalised further by \citet{aggarwal2015generalized} and
given an explicit andness and orness parameterisation by \citet{dujmovic2015weighted}. It is
the natural aggregator when the degree of permitted compensation between criteria is itself
a design question rather than an assumption.

CRITIC \citep{diakoulaki1995determining} derives weights from contrast intensity and
conflict. Objective weighting methods are not interchangeable
\citep{mukhametzyanov2021specific,ariyanti2025comparative}, their stability under
perturbation varies \citep{linh2025evaluating,paradowski2025novel}, they are susceptible to
rank reversal \citep{li2026trend}, and sensitivity analysis over the weight vector is
recommended practice \citep{gaona2025integrated}. We therefore report three schemes and a
rank-stability analysis rather than a single weight vector.

Because the scores feed a downstream decision, calibration matters as much as ranking
\citep{filho2023classifier}, and the excess risk of acting on an uncalibrated classifier is
quantifiable \citep{perezlebel2025decision}. Explanation is subject to its own limits
\citep{bilodeau2024impossibility,hwang2025shapbased}, and in process monitoring specifically
the guidance is to obtain interpretable models rather than to explain opaque ones after the
fact \citep{stevens2024explainability,mehdiyev2025interpretable,sahatova2026feature}.
Whether explanations actually improve human decisions is not settled
\citep{haag2026explanations,chae2025explanation}.
